\chapter{Boolean Analysis — Bool Fourier}
%%% generated by the blueprint pipeline from TCSlib.BooleanAnalysis.BLR.BoolFourier
\section{Overview}

This module develops Fourier analysis on the Boolean hypercube $\{0,1\}^n$, building on
top of \texttt{TCSlib.BooleanAnalysis.Basic}.  It introduces type aliases for the
hypercube and Boolean functions, the group structure (XOR), the uniform expectation
operator, the inner product and $L^2$ norm, Fourier–Walsh characters, Fourier
coefficients, and proves the fundamental theorems: orthonormality of characters, the
Walsh expansion, Parseval's identity, and the convolution theorem.

%%%%%%%%%%%%%%%%%%%%%%%%%%%%%%%%%%%%%%%%%%%%%%%%%%%%%%%%%%%%%%%%%%%%%%%%%%%%%%%
\section{Declarations}

\begin{definition}[Hypercube]
\lean{BoolFourier.hypercube}
\label{BoolFourier.hypercube}
\leanok
\uses{BooleanAnalysis.BoolCube}
The Boolean hypercube of dimension $n$ is the type $\{0,1\}^n$, defined as an
abbreviation for \texttt{BooleanAnalysis.BoolCube} $n$.
\end{definition}

\begin{definition}[Boolean function type]
\lean{BoolFourier.BoolFun}
\label{BoolFourier.BoolFun}
\leanok
\uses{BooleanAnalysis.BooleanFunc}
A Boolean function of arity $n$ is an element of type $\{0,1\}^n \to \mathbb{R}$,
defined as an abbreviation for \texttt{BooleanAnalysis.BooleanFunc} $n$.
\end{definition}

\begin{definition}[Bool-to-$\pm 1$ embedding]
\lean{BoolFourier.BoolToPM1}
\label{BoolFourier.BoolToPM1}
\leanok
\uses{BooleanAnalysis.boolToSign}
The map $\mathrm{BoolToPM1} : \mathrm{Bool} \to \mathbb{R}$ sends $\mathtt{false}$
to $1$ and $\mathtt{true}$ to $-1$, realising the standard identification
$\{0,1\} \cong \{\pm 1\}$.
\end{definition}

\begin{definition}[Partial inverse of $\pm 1$ embedding]
\lean{BoolFourier.PM1ToBool?}
\label{BoolFourier.PM1ToBool?}
\leanok
The partial inverse $\mathrm{PM1ToBool?} : \mathbb{R} \to \mathrm{Option}\,\mathrm{Bool}$
returns $\mathtt{some}\,\mathtt{true}$ on $-1$, $\mathtt{some}\,\mathtt{false}$ on $1$,
and $\mathtt{none}$ on all other real numbers.
\end{definition}

\begin{definition}[All-zeros vector]
\lean{BoolFourier.zero_vec}
\label{BoolFourier.zero_vec}
\leanok
\uses{BoolFourier.hypercube}
$\mathrm{zero\_vec}(n) : \{0,1\}^n$ is the constant-$\mathtt{false}$ function, serving
as the additive identity of the group $(\mathbb{F}_2)^n$.
\end{definition}

\begin{definition}[Componentwise XOR]
\lean{BoolFourier.xor_vec}
\label{BoolFourier.xor_vec}
\leanok
\uses{BoolFourier.hypercube}
For $x, y \in \{0,1\}^n$, the vector $\mathrm{xor\_vec}(x,y)_i = x_i \oplus y_i$ is
the componentwise XOR, giving the group operation on $(\mathbb{F}_2)^n$.
\end{definition}

\begin{lemma}[Cardinality of the Boolean hypercube]
\lean{BoolFourier.card_hypercube}
\label{BoolFourier.card_hypercube}
\leanok
\uses{BoolFourier.hypercube, BooleanAnalysis.BoolCube}
\difficulty{1}
For every natural number $n$, the Boolean hypercube $\{0,1\}^n$ has exactly $2^n$
elements.
\end{lemma}

\begin{lemma}[Mean of the Bool-to-$\pm 1$ embedding vanishes]
\lean{BoolFourier.avg_BoolToPM1}
\label{BoolFourier.avg_BoolToPM1}
\leanok
\uses{BoolFourier.BoolToPM1, BooleanAnalysis.boolToSign}
\difficulty{1}
Under the embedding $\mathrm{BoolToPM1} : \mathrm{Bool} \to \bbr$ that sends
$\mathtt{false}$ to $1$ and $\mathtt{true}$ to $-1$, the average of its two values is
zero:
\[
  \frac{\mathrm{BoolToPM1}(\mathtt{false}) + \mathrm{BoolToPM1}(\mathtt{true})}{2} = 0.
\]
\end{lemma}

\begin{lemma}[Square of the Bool-to-$\pm 1$ embedding]
\lean{BoolFourier.BoolToPM1_sq}
\label{BoolFourier.BoolToPM1_sq}
\leanok
\uses{BoolFourier.BoolToPM1, BooleanAnalysis.boolToSign_mul_self}
\difficulty{1}
For every Boolean value $b$, the image $\mathrm{BoolToPM1}(b)$ under the map sending
$\mathtt{false}$ to $1$ and $\mathtt{true}$ to $-1$ satisfies
$\mathrm{BoolToPM1}(b)\cdot\mathrm{BoolToPM1}(b) = 1$; that is, this $\pm 1$-valued
quantity squares to $1$.
\end{lemma}

\begin{lemma}[Negation under the $\pm 1$ embedding]
\lean{BoolFourier.BoolToPM1_not}
\label{BoolFourier.BoolToPM1_not}
\leanok
\uses{BoolFourier.BoolToPM1, BooleanAnalysis.boolToSign_not}
\difficulty{1}
Let $\mathrm{BoolToPM1} : \mathrm{Bool} \to \bbr$ be the Bool-to-$\pm 1$ embedding,
which sends $\mathtt{false}$ to $1$ and $\mathtt{true}$ to $-1$. Then for every Boolean
value $b$, negating $b$ flips the sign of its image:
\[
  \mathrm{BoolToPM1}(\lnot b) = -\,\mathrm{BoolToPM1}(b).
\]
\end{lemma}

\begin{lemma}[XOR as multiplication under the $\pm 1$ embedding]
\lean{BoolFourier.BoolToPM1_xor}
\label{BoolFourier.BoolToPM1_xor}
\leanok
\uses{BoolFourier.BoolToPM1, BooleanAnalysis.boolToSign}
\difficulty{2}
Let $\mathrm{BoolToPM1} : \mathrm{Bool} \to \bbr$ be the map sending $\mathtt{false}$ to
$1$ and $\mathtt{true}$ to $-1$. Then for all Booleans $a$ and $b$,
\[
  \mathrm{BoolToPM1}(a \oplus b) = \mathrm{BoolToPM1}(a)\cdot\mathrm{BoolToPM1}(b),
\]
so the exclusive-or $a \oplus b$ is carried to the ordinary product in
$\{-1,1\}\subset\bbr$.
\end{lemma}

\begin{definition}[Uniform expectation]
\lean{BoolFourier.expectation}
\label{BoolFourier.expectation}
\leanok
\uses{BoolFourier.BoolFun}
For $f : \{0,1\}^n \to \mathbb{R}$, the uniform expectation is
\[
  \mathbb{E}[f] = \frac{1}{2^n} \sum_{x \in \{0,1\}^n} f(x).
\]
\end{definition}

\begin{lemma}[Factorisation of expectation over product functions]
\lean{BoolFourier.expectation_factorizes}
\label{BoolFourier.expectation_factorizes}
\leanok
\uses{BoolFourier.expectation, BoolFourier.hypercube}
\difficulty{6}
Let $g : \{0,1\}^n \to \bbr$ be given coordinatewise, so that for each index $i \in
\{0,\dots,n-1\}$ we have a function $g_i : \{0,1\} \to \bbr$, and let $f : \{0,1\}^n \to
\bbr$ be the product function $f(x) = \prod_{i} g_i(x_i)$. Then its uniform expectation
over the Boolean hypercube factorises as
\[
  \E[f] = \prod_{i} \frac{g_i(0) + g_i(1)}{2}.
\]
\end{lemma}

\begin{definition}[Inner product of Boolean functions]
\lean{BoolFourier.inner_product}
\label{BoolFourier.inner_product}
\leanok
\uses{BoolFourier.BoolFun, BoolFourier.expectation}
The inner product of $f, g : \{0,1\}^n \to \mathbb{R}$ is
\[
  \langle f, g \rangle = \mathbb{E}[f \cdot g]
  = \frac{1}{2^n}\sum_{x} f(x)\,g(x).
\]
\end{definition}

\begin{definition}[$L^2$ norm squared]
\lean{BoolFourier.L2_norm_sq}
\label{BoolFourier.L2_norm_sq}
\leanok
\uses{BoolFourier.BoolFun, BoolFourier.expectation}
The squared $L^2$ norm of $f : \{0,1\}^n \to \mathbb{R}$ is
\[
  \|f\|_2^2 = \mathbb{E}[f^2]
  = \frac{1}{2^n}\sum_{x} f(x)^2.
\]
\end{definition}

\begin{definition}[Convolution on the hypercube]
\lean{BoolFourier.convolution}
\label{BoolFourier.convolution}
\leanok
\statementsource{boolean-ch08-generalized-domains}{pdf-50833afe7403-p023-b004}
\uses{BoolFourier.BoolFun, BoolFourier.expectation, BoolFourier.xor_vec}
The convolution of $f, g : \{0,1\}^n \to \mathbb{R}$ is the function
\[
  (f * g)(x) = \mathbb{E}_{y}\bigl[f(y)\,g(x \oplus y)\bigr].
\]
\end{definition}

\begin{definition}[Walsh--Fourier character]
\lean{BoolFourier.char_S}
\label{BoolFourier.char_S}
\leanok
\statementsource{boolean-ch08-generalized-domains}{pdf-50833afe7403-p002-b007}
\uses{BoolFourier.BoolFun, BooleanAnalysis.chiS}
For a set $S \subseteq [n]$, the Walsh character
$\chi_S : \{0,1\}^n \to \mathbb{R}$ is
\[
  \chi_S(x) = \prod_{i \in S} (-1)^{x_i},
\]
defined as an alias for \texttt{BooleanAnalysis.chiS}.
\end{definition}

\begin{lemma}[Walsh character at the origin]
\lean{BoolFourier.char_S_of_zero}
\label{BoolFourier.char_S_of_zero}
\leanok
\uses{BoolFourier.char_S, BoolFourier.zero_vec, BooleanAnalysis.boolToSign, BooleanAnalysis.chiS}
\difficulty{1}
Let $n$ be a natural number and let $S \subseteq [n]$. Evaluating the Walsh character
$\chi_S : \{0,1\}^n \to \bbr$ at the all-zeros point $\mathbf{0}$ (the vector whose
every coordinate is $\mathtt{false}$) gives
\[
  \chi_S(\mathbf{0}) = 1.
\]
\end{lemma}

\begin{lemma}[Walsh characters take values in the unit sphere]
\lean{BoolFourier.char_S_sq}
\label{BoolFourier.char_S_sq}
\leanok
\uses{BoolFourier.char_S, BoolFourier.hypercube, BooleanAnalysis.chiS_sq_eq_one}
\difficulty{2}
Let $n$ be a natural number, let $S \subseteq [n]$, and let $x \in \{0,1\}^n$. Then the
Walsh–Fourier character $\chi_S(x) = \prod_{i \in S}(-1)^{x_i}$ satisfies
\[
  \chi_S(x)\,\chi_S(x) = 1.
\]
\end{lemma}

\begin{lemma}[Empty Walsh character is the constant one]
\lean{BoolFourier.char_S_empty}
\label{BoolFourier.char_S_empty}
\leanok
\uses{BoolFourier.char_S, BooleanAnalysis.chiS_empty}
\difficulty{1}
Let $n$ be a natural number. The Walsh character $\chi_\emptyset$ associated to the
empty subset of $[n]$ is the constant function equal to $1$ on the Boolean hypercube;
that is, $\chi_\emptyset(x) = 1$ for every $x \in \{0,1\}^n$.
\end{lemma}

\begin{lemma}[Product of Walsh characters]
\lean{BoolFourier.char_S_times_char_T}
\proofsource{boolean-ch08-generalized-domains}{pdf-50833afe7403-p023-b002}
\label{BoolFourier.char_S_times_char_T}
\leanok
\uses{BoolFourier.char_S, BoolFourier.hypercube, BooleanAnalysis.chiS_mul_chiS}
\difficulty{1}
For any two subsets $S, T \subseteq [n]$ and any point $x \in \{0,1\}^n$,
\[
  \chi_S(x)\,\chi_T(x) = \chi_{S \mathbin{\triangle} T}(x),
\]
where $\chi_S$ denotes the Walsh--Fourier character of $S$ and $S \mathbin{\triangle} T$
is the symmetric difference of $S$ and $T$.
\end{lemma}

\begin{lemma}[Summing all Walsh characters at the origin]
\lean{BoolFourier.sum_char_S_at_zero}
\label{BoolFourier.sum_char_S_at_zero}
\leanok
\uses{BoolFourier.char_S, BoolFourier.char_S_of_zero, BoolFourier.hypercube, BoolFourier.zero_vec}
\difficulty{1}
Fix $n \in \bbn$, and for each subset $S \subseteq [n]$ let $\chi_S \colon \{0,1\}^n \to
\bbr$ denote the Walsh--Fourier character $\chi_S(x) = \prod_{i \in S}(-1)^{x_i}$.
Writing $\mathbf{0}$ for the all-zeros point of $\{0,1\}^n$, summing these characters
evaluated at $\mathbf{0}$ over all $2^n$ subsets of $[n]$ gives
\[
  \sum_{S \subseteq [n]} \chi_S(\mathbf{0}) \;=\; \abs{\{0,1\}^n},
\]
the number of points of the Boolean hypercube.
\end{lemma}

\begin{lemma}[Vanishing of the character sum at a nonzero point]
\lean{BoolFourier.sum_char_S_ne_zero}
\label{BoolFourier.sum_char_S_ne_zero}
\leanok
\uses{BoolFourier.char_S, BoolFourier.hypercube, BoolFourier.zero_vec, BooleanAnalysis.boolToSign, BooleanAnalysis.chiS}
\difficulty{4}
Fix $n\in\bbn$ and a point $x\in\{0,1\}^n$, and for each subset $S\subseteq[n]$ let
$\chi_S(x)=\prod_{i\in S}(-1)^{x_i}$ be the Walsh–Fourier character. If $x$ is not the
all-zeros point, then
\[
  \sum_{S\subseteq[n]}\chi_S(x)=0,
\]
the sum ranging over all $2^n$ subsets of $[n]$.
\end{lemma}

\begin{lemma}[Expectation of the empty Walsh character]
\lean{BoolFourier.expectation_char_empty}
\label{BoolFourier.expectation_char_empty}
\leanok
\uses{BoolFourier.char_S, BoolFourier.expectation}
\difficulty{1}
For every $n \in \bbn$, the Walsh character $\chi_\emptyset$ associated to the empty
subset of $[n]$ has uniform expectation equal to one:
\[
\E[\chi_\emptyset] \;=\; \frac{1}{2^n}\sum_{x \in \{0,1\}^n} \chi_\emptyset(x) \;=\; 1.
\]
\end{lemma}

\begin{lemma}[Non-trivial characters have mean zero]
\lean{BoolFourier.expectation_char_nonempty}
\statementsource{boolean-ch08-generalized-domains}{
  pdf-50833afe7403-p022-b002,
  pdf-50833afe7403-p022-b003
}
\label{BoolFourier.expectation_char_nonempty}
\leanok
\proofsource{boolean-ch08-generalized-domains}{
  pdf-50833afe7403-p004-b006,
  pdf-50833afe7403-p003-b003
}
\uses{BoolFourier.BoolToPM1, BoolFourier.char_S, BoolFourier.expectation, BoolFourier.expectation_factorizes}
\difficulty{1}
Let $n$ be a natural number and let $S \subseteq [n]$ be a nonempty set of coordinates.
Then the Walsh--Fourier character $\chi_S(x) = \prod_{i \in S}(-1)^{x_i}$ has vanishing
uniform expectation over the Boolean hypercube $\{0,1\}^n$:
\[
  \E[\chi_S] = \frac{1}{2^n}\sum_{x \in \{0,1\}^n} \chi_S(x) = 0.
\]
\end{lemma}

\begin{lemma}[Self inner product of a Walsh character]
\lean{BoolFourier.inner_product_char_self}
\statementsource{boolean-ch08-generalized-domains}{
  pdf-50833afe7403-p004-b006,
  pdf-50833afe7403-p004-b005
}
\label{BoolFourier.inner_product_char_self}
\leanok
\proofsource{boolean-ch08-generalized-domains}{
  pdf-50833afe7403-p004-b006,
  pdf-50833afe7403-p004-b005
}
\uses{BoolFourier.char_S, BoolFourier.char_S_sq, BoolFourier.expectation, BoolFourier.inner_product}
\difficulty{3}
Let $n$ be a natural number and let $S \subseteq [n]$. Then the Walsh--Fourier character
$\chi_S(x) = \prod_{i \in S}(-1)^{x_i}$ satisfies
\[
\langle \chi_S, \chi_S \rangle \;=\; \frac{1}{2^n}\sum_{x \in \{0,1\}^n} \chi_S(x)^2
\;=\; 1,
\]
where the inner product is the uniform expectation of the pointwise product.
\end{lemma}

\begin{lemma}[Orthogonality of distinct Walsh characters]
\lean{BoolFourier.inner_product_char_nonself}
\label{BoolFourier.inner_product_char_nonself}
\leanok
\proofsource{boolean-ch08-generalized-domains}{
  pdf-50833afe7403-p004-b006,
  pdf-50833afe7403-p004-b005
}
\proofstep
  {BoolCircuit.Lit}
  {context}
  {boolean-ch08-generalized-domains}
  {pdf-50833afe7403-p032-b002}
\proofstep
  {BoolCircuit.Lit.eval}
  {context}
  {boolean-ch08-generalized-domains}
  {pdf-50833afe7403-p032-b002}
\proofstep
  {Literal}
  {context}
  {boolean-ch08-generalized-domains}
  {pdf-50833afe7403-p032-b002}
\proofstep
  {DecisionTree.eval}
  {context}
  {boolean-ch08-generalized-domains}
  {pdf-50833afe7403-p024-b001}
\proofstep
  {DecisionTree}
  {context}
  {boolean-ch08-generalized-domains}
  {pdf-50833afe7403-p023-b009}
\proofstep
  {DecisionTree.depth}
  {context}
  {boolean-ch08-generalized-domains}
  {pdf-50833afe7403-p024-b001}
\proofstep
  {SwitchingLemma2.Restriction}
  {context}
  {boolean-ch08-generalized-domains}
  {pdf-50833afe7403-p002-b003}
\proofstep
  {SwitchingLemma2.Literal.fixedBy}
  {context}
  {boolean-ch08-generalized-domains}
  {pdf-50833afe7403-p002-b003}
\proofstep
  {Literal.flipNeg}
  {context}
  {boolean-ch08-generalized-domains}
  {pdf-50833afe7403-p032-b002}
\proofstep
  {Literal.flipNeg_eval}
  {context}
  {boolean-ch08-generalized-domains}
  {pdf-50833afe7403-p032-b002}
\proofstep
  {Literal.flipNeg_var}
  {context}
  {boolean-ch08-generalized-domains}
  {pdf-50833afe7403-p032-b002}
\proofstep
  {DecisionTree.negateLeaves}
  {context}
  {boolean-ch08-generalized-domains}
  {pdf-50833afe7403-p024-b001}
\proofstep
  {DecisionTree.negateLeaves_eval}
  {context}
  {boolean-ch08-generalized-domains}
  {pdf-50833afe7403-p024-b001}
\proofstep
  {DecisionTree.negateLeaves_depth}
  {context}
  {boolean-ch08-generalized-domains}
  {pdf-50833afe7403-p024-b001}
\proofstep
  {LMN.clauseIsTaut}
  {context}
  {boolean-ch08-generalized-domains}
  {pdf-50833afe7403-p032-b002}
\proofstep
  {LMN.instDecidableClauseIsTaut}
  {context}
  {boolean-ch08-generalized-domains}
  {pdf-50833afe7403-p032-b002}
\proofstep
  {LMN.mergeGates}
  {context}
  {boolean-ch08-generalized-domains}
  {pdf-50833afe7403-p039-b001}
\proofstep
  {LMN.mergeGates_castAdd}
  {context}
  {boolean-ch08-generalized-domains}
  {pdf-50833afe7403-p039-b001}
\proofstep
  {LMN.mergeGates_natAdd}
  {context}
  {boolean-ch08-generalized-domains}
  {pdf-50833afe7403-p039-b001}
\proofstep
  {BooleanAnalysis.BoolCube}
  {context}
  {boolean-ch08-generalized-domains}
  {pdf-50833afe7403-p002-b002}
\proofstep
  {BooleanAnalysis.BooleanFunc}
  {context}
  {boolean-ch08-generalized-domains}
  {pdf-50833afe7403-p001-b004}
\proofstep
  {BooleanAnalysis.uniformWeight}
  {context}
  {boolean-ch08-generalized-domains}
  {pdf-50833afe7403-p002-b002}
\proofstep
  {BooleanAnalysis.expect}
  {context}
  {boolean-ch08-generalized-domains}
  {pdf-50833afe7403-p001-b004}
\proofstep
  {BooleanAnalysis.innerProduct}
  {direct}
  {boolean-ch08-generalized-domains}
  {pdf-50833afe7403-p001-b004}
\proofstep
  {BooleanAnalysis.boolToSign}
  {context}
  {boolean-ch08-generalized-domains}
  {pdf-50833afe7403-p002-b007}
\proofstep
  {BooleanAnalysis.boolToSign_false}
  {context}
  {boolean-ch08-generalized-domains}
  {pdf-50833afe7403-p002-b007}
\proofstep
  {BooleanAnalysis.boolToSign_true}
  {context}
  {boolean-ch08-generalized-domains}
  {pdf-50833afe7403-p002-b007}
\proofstep
  {BooleanAnalysis.boolToSign_sq}
  {context}
  {boolean-ch08-generalized-domains}
  {pdf-50833afe7403-p003-b002}
\proofstep
  {BooleanAnalysis.boolToSign_mul_self}
  {context}
  {boolean-ch08-generalized-domains}
  {pdf-50833afe7403-p003-b002}
\proofstep
  {BooleanAnalysis.chiS}
  {direct}
  {boolean-ch08-generalized-domains}
  {pdf-50833afe7403-p002-b007}
\proofstep
  {BooleanAnalysis.chiS_empty}
  {context}
  {boolean-ch08-generalized-domains}
  {pdf-50833afe7403-p003-b003}
\proofstep
  {BooleanAnalysis.chiS_singleton}
  {direct}
  {boolean-ch08-generalized-domains}
  {pdf-50833afe7403-p002-b007}
\proofstep
  {BooleanAnalysis.chiS_mul_chiS}
  {direct}
  {boolean-ch08-generalized-domains}
  {pdf-50833afe7403-p023-b002,
    pdf-50833afe7403-p015-b002}
\proofstep
  {BooleanAnalysis.sum_boolToSign}
  {direct}
  {boolean-ch08-generalized-domains}
  {pdf-50833afe7403-p003-b002}
\proofstep
  {BooleanAnalysis.sum_chiS}
  {direct}
  {boolean-ch08-generalized-domains}
  {pdf-50833afe7403-p022-b002,
    pdf-50833afe7403-p003-b003}
\proofstep
  {BooleanAnalysis.fourier_coeff_chi}
  {direct}
  {boolean-ch08-generalized-domains}
  {pdf-50833afe7403-p004-b006,
    pdf-50833afe7403-p003-b002}
\proofstep
  {BooleanAnalysis.innerProduct_chi_self}
  {context}
  {boolean-ch08-generalized-domains}
  {pdf-50833afe7403-p003-b002}
\proofstep
  {BooleanAnalysis.flipBit}
  {context}
  {boolean-ch08-generalized-domains}
  {pdf-50833afe7403-p017-b001}
\proofstep
  {BooleanAnalysis.flipBit_flipBit}
  {context}
  {boolean-ch08-generalized-domains}
  {pdf-50833afe7403-p017-b001}
\proofstep
  {BooleanAnalysis.instAddCommGroupBooleanFunc}
  {context}
  {boolean-ch08-generalized-domains}
  {pdf-50833afe7403-p001-b004}
\proofstep
  {BooleanAnalysis.instModuleRealBooleanFunc}
  {context}
  {boolean-ch08-generalized-domains}
  {pdf-50833afe7403-p001-b004}
\proofstep
  {BooleanAnalysis.boolToSign_not}
  {context}
  {boolean-ch08-generalized-domains}
  {pdf-50833afe7403-p002-b007}
\proofstep
  {BoolFourier.hypercube}
  {context}
  {boolean-ch08-generalized-domains}
  {pdf-50833afe7403-p002-b002}
\proofstep
  {BoolFourier.BoolFun}
  {context}
  {boolean-ch08-generalized-domains}
  {pdf-50833afe7403-p001-b004}
\proofstep
  {BoolFourier.BoolToPM1}
  {context}
  {boolean-ch08-generalized-domains}
  {pdf-50833afe7403-p002-b007}
\proofstep
  {BoolFourier.expectation}
  {context}
  {boolean-ch08-generalized-domains}
  {pdf-50833afe7403-p001-b004}
\proofstep
  {BoolFourier.expectation_factorizes}
  {context}
  {boolean-ch08-generalized-domains}
  {pdf-50833afe7403-p004-b006}
\proofstep
  {BoolFourier.inner_product}
  {direct}
  {boolean-ch08-generalized-domains}
  {pdf-50833afe7403-p001-b004}
\proofstep
  {BoolFourier.char_S}
  {direct}
  {boolean-ch08-generalized-domains}
  {pdf-50833afe7403-p002-b007}
\proofstep
  {BoolFourier.char_S_times_char_T}
  {direct}
  {boolean-ch08-generalized-domains}
  {pdf-50833afe7403-p023-b002,
    pdf-50833afe7403-p015-b002}
\proofstep
  {BoolFourier.expectation_char_nonempty_aux_h_exp}
  {context}
  {boolean-ch08-generalized-domains}
  {pdf-50833afe7403-p004-b006}
\proofstep
  {BoolFourier.expectation_char_nonempty}
  {direct}
  {boolean-ch08-generalized-domains}
  {pdf-50833afe7403-p022-b002,
    pdf-50833afe7403-p003-b003}
\uses{BoolFourier.char_S, BoolFourier.char_S_times_char_T, BoolFourier.expectation_char_nonempty, BoolFourier.inner_product}
\difficulty{3}
Fix $n \in \bbn$, and let $S, T \subseteq [n]$ be two distinct subsets. Then the
Walsh--Fourier characters $\chi_S$ and $\chi_T$ are orthogonal with respect to the
uniform inner product on $\{0,1\}^n$:
\[
\langle \chi_S, \chi_T \rangle = \frac{1}{2^n}\sum_{x \in \{0,1\}^n}
\chi_S(x)\,\chi_T(x) = 0.
\]
\end{lemma}

\begin{definition}[Fourier coefficient]
\lean{BoolFourier.fourier_coeff}
\label{BoolFourier.fourier_coeff}
\leanok
\statementsource{boolean-ch08-generalized-domains}{pdf-50833afe7403-p005-b001}
\uses{BoolFourier.BoolFun, BoolFourier.char_S, BoolFourier.inner_product}
The Fourier coefficient of $f : \{0,1\}^n \to \mathbb{R}$ at $S \subseteq [n]$ is
\[
  \hat f(S) = \langle f, \chi_S \rangle.
\]
\end{definition}

\begin{lemma}[Self-Fourier coefficient of a Walsh character]
\lean{BoolFourier.fourier_coeff_char_self}
\label{BoolFourier.fourier_coeff_char_self}
\leanok
\proofsource{boolean-ch08-generalized-domains}{pdf-50833afe7403-p004-b006}
\uses{BoolFourier.char_S, BoolFourier.fourier_coeff, BoolFourier.inner_product_char_self}
\difficulty{1}
For every subset $S \subseteq [n]$, the Fourier coefficient of the Walsh character
$\chi_S$ at the same set $S$ equals $1$; that is,
\[
  \widehat{\chi_S}(S) = \langle \chi_S, \chi_S \rangle = 1.
\]
\end{lemma}

\begin{lemma}[Orthogonality of Walsh characters]
\lean{BoolFourier.fourier_coeff_char_of_ne}
\label{BoolFourier.fourier_coeff_char_of_ne}
\leanok
\proofsource{boolean-ch08-generalized-domains}{pdf-50833afe7403-p004-b006}
\uses{BoolFourier.char_S, BoolFourier.fourier_coeff, BoolFourier.inner_product_char_nonself}
\difficulty{1}
Let $S, T \subseteq [n]$ be distinct subsets of coordinates. Then the Fourier
coefficient of the Walsh character $\chi_S : \{0,1\}^n \to \bbr$ at $T$ vanishes:
\[
  \widehat{\chi_S}(T) = \langle \chi_S, \chi_T \rangle = 0.
\]
\end{lemma}

\begin{lemma}[Agreement of two definitions of the uniform expectation]
\lean{BoolFourier.expectation_eq_expect}
\label{BoolFourier.expectation_eq_expect}
\leanok
\uses{BoolFourier.BoolFun, BoolFourier.expectation, BooleanAnalysis.expect, BooleanAnalysis.uniformWeight}
\difficulty{2}
Let $n$ be a natural number and let $f : \{0,1\}^n \to \bbr$ be a Boolean function of
arity $n$. Then the two forms of the uniform expectation of $f$ coincide, namely
\[
  \frac{1}{2^n} \sum_{x \in \{0,1\}^n} f(x) \;=\; 2^{-n} \sum_{x \in \{0,1\}^n} f(x).
\]
\end{lemma}

\begin{lemma}[Agreement of two Fourier coefficient definitions]
\lean{BoolFourier.fourier_coeff_eq}
\label{BoolFourier.fourier_coeff_eq}
\leanok
\uses{BoolFourier.BoolFun, BoolFourier.expectation_eq_expect, BoolFourier.fourier_coeff, BoolFourier.inner_product, BooleanAnalysis.fourierCoeff, BooleanAnalysis.innerProduct}
\difficulty{2}
Let $f : \{0,1\}^n \to \bbr$ be a Boolean function of arity $n$, and let $S \subseteq
[n]$. Then the two definitions of the Fourier–Walsh coefficient $\hat f(S)$ coincide:
both are equal to the $L^2$ inner product of $f$ with the Walsh character $\chi_S$,
\[
\hat f(S) \;=\; \langle f, \chi_S\rangle \;=\; 2^{-n}\sum_{x\in\{0,1\}^n}
f(x)\,\chi_S(x).
\]
\end{lemma}

\begin{lemma}[Fourier–Walsh expansion]
\lean{BoolFourier.fourier_expansion}
\label{BoolFourier.fourier_expansion}
\leanok
\statementsource{boolean-ch08-generalized-domains}{pdf-50833afe7403-p005-b001}
\uses{BoolFourier.BoolFun, BoolFourier.char_S, BoolFourier.fourier_coeff, BoolFourier.fourier_coeff_eq, BoolFourier.hypercube, BooleanAnalysis.walsh_expansion}
\difficulty{2}
Let $f : \{0,1\}^n \to \bbr$ be a Boolean function and let $x \in \{0,1\}^n$. Then $f$
agrees pointwise with the sum of its Walsh characters weighted by its Fourier
coefficients:
\[
  f(x) = \sum_{S \subseteq [n]} \hat f(S)\,\chi_S(x),
\]
where the sum ranges over all subsets $S$ of $[n]$, $\hat f(S) = \langle f,
\chi_S\rangle$ is the Fourier coefficient of $f$ at $S$, and $\chi_S$ is the Walsh
character of $S$.
\end{lemma}

\begin{lemma}[Parseval's identity for the Walsh–Fourier expansion]
\lean{BoolFourier.parseval_identity}
\label{BoolFourier.parseval_identity}
\leanok
\proofsource{boolean-ch08-generalized-domains}{
  pdf-50833afe7403-p005-b006,
  pdf-50833afe7403-p006-b001
}
\proofstep
  {BoolCircuit.Lit}
  {context}
  {boolean-ch08-generalized-domains}
  {pdf-50833afe7403-p023-b009}
\proofstep
  {BoolCircuit.Lit.eval}
  {context}
  {boolean-ch08-generalized-domains}
  {pdf-50833afe7403-p023-b009}
\proofstep
  {Literal}
  {context}
  {boolean-ch08-generalized-domains}
  {pdf-50833afe7403-p032-b002}
\proofstep
  {DecisionTree.eval}
  {context}
  {boolean-ch08-generalized-domains}
  {pdf-50833afe7403-p023-b009}
\proofstep
  {DecisionTree}
  {context}
  {boolean-ch08-generalized-domains}
  {pdf-50833afe7403-p023-b009}
\proofstep
  {DecisionTree.depth}
  {context}
  {boolean-ch08-generalized-domains}
  {pdf-50833afe7403-p024-b001}
\proofstep
  {SwitchingLemma2.Restriction}
  {context}
  {boolean-ch08-generalized-domains}
  {pdf-50833afe7403-p027-b002}
\proofstep
  {SwitchingLemma2.Literal.fixedBy}
  {context}
  {boolean-ch08-generalized-domains}
  {pdf-50833afe7403-p027-b002}
\proofstep
  {Literal.flipNeg}
  {context}
  {boolean-ch08-generalized-domains}
  {pdf-50833afe7403-p023-b009}
\proofstep
  {Literal.flipNeg_eval}
  {context}
  {boolean-ch08-generalized-domains}
  {pdf-50833afe7403-p023-b009}
\proofstep
  {Literal.flipNeg_var}
  {context}
  {boolean-ch08-generalized-domains}
  {pdf-50833afe7403-p023-b009}
\proofstep
  {DecisionTree.negateLeaves}
  {context}
  {boolean-ch08-generalized-domains}
  {pdf-50833afe7403-p023-b009}
\proofstep
  {DecisionTree.negateLeaves_eval}
  {context}
  {boolean-ch08-generalized-domains}
  {pdf-50833afe7403-p023-b009}
\proofstep
  {DecisionTree.negateLeaves_depth}
  {context}
  {boolean-ch08-generalized-domains}
  {pdf-50833afe7403-p024-b001}
\proofstep
  {LMN.clauseIsTaut}
  {context}
  {boolean-ch08-generalized-domains}
  {pdf-50833afe7403-p032-b002}
\proofstep
  {LMN.instDecidableClauseIsTaut}
  {context}
  {boolean-ch08-generalized-domains}
  {pdf-50833afe7403-p032-b002}
\proofstep
  {LMN.mergeGates}
  {context}
  {boolean-ch08-generalized-domains}
  {pdf-50833afe7403-p039-b001}
\proofstep
  {LMN.mergeGates_castAdd}
  {context}
  {boolean-ch08-generalized-domains}
  {pdf-50833afe7403-p039-b001}
\proofstep
  {LMN.mergeGates_natAdd}
  {context}
  {boolean-ch08-generalized-domains}
  {pdf-50833afe7403-p039-b001}
\proofstep
  {BooleanAnalysis.BoolCube}
  {context}
  {boolean-ch08-generalized-domains}
  {pdf-50833afe7403-p002-b002}
\proofstep
  {BooleanAnalysis.BooleanFunc}
  {context}
  {boolean-ch08-generalized-domains}
  {pdf-50833afe7403-p001-b004}
\proofstep
  {BooleanAnalysis.uniformWeight}
  {context}
  {boolean-ch08-generalized-domains}
  {pdf-50833afe7403-p001-b004}
\proofstep
  {BooleanAnalysis.expect}
  {context}
  {boolean-ch08-generalized-domains}
  {pdf-50833afe7403-p001-b004}
\proofstep
  {BooleanAnalysis.innerProduct}
  {direct}
  {boolean-ch08-generalized-domains}
  {pdf-50833afe7403-p001-b004}
\proofstep
  {BooleanAnalysis.boolToSign}
  {context}
  {boolean-ch08-generalized-domains}
  {pdf-50833afe7403-p002-b002}
\proofstep
  {BooleanAnalysis.boolToSign_false}
  {context}
  {boolean-ch08-generalized-domains}
  {pdf-50833afe7403-p002-b002}
\proofstep
  {BooleanAnalysis.boolToSign_true}
  {context}
  {boolean-ch08-generalized-domains}
  {pdf-50833afe7403-p002-b002}
\proofstep
  {BooleanAnalysis.boolToSign_sq}
  {context}
  {boolean-ch08-generalized-domains}
  {pdf-50833afe7403-p003-b002}
\proofstep
  {BooleanAnalysis.boolToSign_mul_self}
  {context}
  {boolean-ch08-generalized-domains}
  {pdf-50833afe7403-p003-b002}
\proofstep
  {BooleanAnalysis.chiS}
  {direct}
  {boolean-ch08-generalized-domains}
  {pdf-50833afe7403-p002-b007}
\proofstep
  {BooleanAnalysis.chiS_empty}
  {context}
  {boolean-ch08-generalized-domains}
  {pdf-50833afe7403-p003-b003}
\proofstep
  {BooleanAnalysis.chiS_singleton}
  {context}
  {boolean-ch08-generalized-domains}
  {pdf-50833afe7403-p002-b007}
\proofstep
  {BooleanAnalysis.chiS_mul_chiS}
  {context}
  {boolean-ch08-generalized-domains}
  {pdf-50833afe7403-p015-b002}
\proofstep
  {BooleanAnalysis.fourierCoeff}
  {direct}
  {boolean-ch08-generalized-domains}
  {pdf-50833afe7403-p005-b001}
\proofstep
  {BooleanAnalysis.sum_prod_subset_eq_prod_one_add}
  {context}
  {boolean-ch08-generalized-domains}
  {pdf-50833afe7403-p005-b001}
\proofstep
  {BooleanAnalysis.sum_chiS_mul_eq}
  {context}
  {boolean-ch08-generalized-domains}
  {pdf-50833afe7403-p005-b001}
\proofstep
  {BooleanAnalysis.walsh_expansion}
  {direct}
  {boolean-ch08-generalized-domains}
  {pdf-50833afe7403-p005-b001}
\proofstep
  {BooleanAnalysis.sum_boolToSign}
  {context}
  {boolean-ch08-generalized-domains}
  {pdf-50833afe7403-p003-b002}
\proofstep
  {BooleanAnalysis.sum_chiS}
  {context}
  {boolean-ch08-generalized-domains}
  {pdf-50833afe7403-p003-b003}
\proofstep
  {BooleanAnalysis.fourier_coeff_chi}
  {direct}
  {boolean-ch08-generalized-domains}
  {pdf-50833afe7403-p004-b006}
\proofstep
  {BooleanAnalysis.innerProduct_chi_self}
  {context}
  {boolean-ch08-generalized-domains}
  {pdf-50833afe7403-p004-b006}
\proofstep
  {BooleanAnalysis.parseval}
  {direct}
  {boolean-ch08-generalized-domains}
  {pdf-50833afe7403-p005-b006}
\proofstep
  {BooleanAnalysis.flipBit}
  {context}
  {boolean-ch08-generalized-domains}
  {pdf-50833afe7403-p017-b001}
\proofstep
  {BooleanAnalysis.flipBit_flipBit}
  {context}
  {boolean-ch08-generalized-domains}
  {pdf-50833afe7403-p017-b001}
\proofstep
  {BooleanAnalysis.derivative}
  {context}
  {boolean-ch08-generalized-domains}
  {pdf-50833afe7403-p016-b004}
\proofstep
  {BooleanAnalysis.derivative_add}
  {context}
  {boolean-ch08-generalized-domains}
  {pdf-50833afe7403-p016-b004}
\proofstep
  {BooleanAnalysis.derivative_smul}
  {context}
  {boolean-ch08-generalized-domains}
  {pdf-50833afe7403-p016-b004}
\proofstep
  {BooleanAnalysis.derivativeLm}
  {context}
  {boolean-ch08-generalized-domains}
  {pdf-50833afe7403-p016-b004}
\proofstep
  {BooleanAnalysis.instAddCommGroupBooleanFunc}
  {context}
  {boolean-ch08-generalized-domains}
  {pdf-50833afe7403-p001-b004}
\proofstep
  {BooleanAnalysis.instModuleRealBooleanFunc}
  {context}
  {boolean-ch08-generalized-domains}
  {pdf-50833afe7403-p001-b004}
\proofstep
  {BooleanAnalysis.boolToSign_not}
  {context}
  {boolean-ch08-generalized-domains}
  {pdf-50833afe7403-p002-b002}
\proofstep
  {BoolFourier.BoolFun}
  {context}
  {boolean-ch08-generalized-domains}
  {pdf-50833afe7403-p001-b004}
\proofstep
  {BoolFourier.expectation}
  {context}
  {boolean-ch08-generalized-domains}
  {pdf-50833afe7403-p001-b004}
\proofstep
  {BoolFourier.inner_product}
  {direct}
  {boolean-ch08-generalized-domains}
  {pdf-50833afe7403-p001-b004}
\proofstep
  {BoolFourier.L2_norm_sq}
  {context}
  {boolean-ch08-generalized-domains}
  {pdf-50833afe7403-p002-b003}
\proofstep
  {BoolFourier.char_S}
  {direct}
  {boolean-ch08-generalized-domains}
  {pdf-50833afe7403-p002-b007}
\proofstep
  {BoolFourier.fourier_coeff}
  {direct}
  {boolean-ch08-generalized-domains}
  {pdf-50833afe7403-p005-b001}
\proofstep
  {BoolFourier.expectation_eq_expect}
  {context}
  {boolean-ch08-generalized-domains}
  {pdf-50833afe7403-p001-b004}
\proofstep
  {BoolFourier.fourier_coeff_eq}
  {context}
  {boolean-ch08-generalized-domains}
  {pdf-50833afe7403-p005-b001}
\proofstep
  {BoolFourier.parseval_identity_aux_hL2}
  {context}
  {boolean-ch08-generalized-domains}
  {pdf-50833afe7403-p005-b006}
\uses{BoolFourier.BoolFun, BoolFourier.L2_norm_sq, BoolFourier.expectation_eq_expect, BoolFourier.fourier_coeff, BoolFourier.fourier_coeff_eq, BooleanAnalysis.innerProduct, BooleanAnalysis.parseval}
\difficulty{2}
Let $f : \{0,1\}^n \to \bbr$ be a real-valued function on the Boolean hypercube, and for
each subset $S \subseteq [n]$ let $\hat f(S) = \langle f, \chi_S\rangle$ be its Fourier
coefficient, where the inner product and expectation are taken with respect to the
uniform measure. Then the sum of the squared Fourier coefficients equals the squared
$L^2$ norm of $f$:
\[
  \sum_{S \subseteq [n]} \hat f(S)^2 \;=\; \|f\|_2^2 \;=\; \E\bigl[f^2\bigr].
\]
\end{lemma}

\begin{lemma}[Convolution theorem on the Boolean cube]
\lean{BoolFourier.fourier_coeff_convolution}
\statementsource{boolean-ch08-generalized-domains}{pdf-50833afe7403-p023-b006}
\label{BoolFourier.fourier_coeff_convolution}
\leanok
\uses{BoolFourier.BoolFun, BoolFourier.char_S, BoolFourier.convolution, BoolFourier.expectation, BoolFourier.fourier_coeff, BoolFourier.hypercube, BoolFourier.inner_product, BoolFourier.xor_vec, BooleanAnalysis.boolToSign, BooleanAnalysis.chiS}
\difficulty{4}
Let $n$ be a natural number, and let $f, g : \{0,1\}^n \to \bbr$ be Boolean functions.
Then for every subset $S \subseteq [n]$, the Fourier coefficient of the convolution $f *
g$ at $S$ equals the product of the Fourier coefficients of $f$ and $g$ at $S$:
\[
  \widehat{f * g}(S) = \hat f(S)\,\hat g(S).
\]
\end{lemma}

%%% added by the blueprint pipeline
\section{Additional declarations}

\begin{lemma}[Walsh characters square to one]
\lean{BoolFourier.char_S_sq_aux_h}
\label{BoolFourier.char_S_sq_aux_h}
\leanok
\uses{BoolFourier.hypercube, BooleanAnalysis.chiS_sq_eq_one}
\difficulty{1}
Fix a natural number $n$, and let $S \subseteq \{1,\dots,n\}$ be a set of coordinates.
For every point $x$ of the Boolean hypercube $\{0,1\}^n$, the Walsh character $\chi_S$
satisfies
\[
\chi_S(x)^2 = 1.
\]
\end{lemma}

\begin{lemma}[Walsh characters equal one at the origin]
\lean{BoolFourier.sum_char_S_at_zero_aux_h}
\label{BoolFourier.sum_char_S_at_zero_aux_h}
\leanok
\uses{BoolFourier.char_S, BoolFourier.char_S_of_zero, BoolFourier.zero_vec}
\difficulty{1}
Fix $n \ge 0$, and for each subset $S \subseteq [n]$ let $\chi_S : \{0,1\}^n \to \bbr$
be the Walsh character $\chi_S(x) = \prod_{i \in S}(-1)^{x_i}$. Writing $\mathbf{0}$ for
the all-zeros point of $\{0,1\}^n$, the map sending each $S$ to $\chi_S(\mathbf{0})$ is
the constant function equal to $1$:
\[
  \bigl(S \mapsto \chi_S(\mathbf{0})\bigr) \;=\; \bigl(S \mapsto 1\bigr).
\]
\end{lemma}

\begin{lemma}[Sign-flip symmetry of the character sum at a nonzero point]
\lean{BoolFourier.sum_char_S_ne_zero_aux_h_pair}
\label{BoolFourier.sum_char_S_ne_zero_aux_h_pair}
\leanok
\uses{BoolFourier.char_S, BoolFourier.hypercube, BoolFourier.zero_vec, BooleanAnalysis.boolToSign, BooleanAnalysis.chiS}
\difficulty{6}
Let $x \in \{0,1\}^n$ be a point of the Boolean hypercube that differs from the
all-zeros point $\mathbf{0}$, and suppose $i$ is a coordinate at which $x_i =
\mathtt{true}$. Then summing the Walsh–Fourier characters over all subsets $S \subseteq
[n]$, evaluated at $x$, gives the same value as summing their negations:
\[
  \sum_{S \subseteq [n]} \chi_S(x) \;=\; \sum_{S \subseteq [n]} \bigl(-\chi_S(x)\bigr).
\]
\end{lemma}

\begin{lemma}[Expectation of a Walsh character as a coordinate product]
\lean{BoolFourier.expectation_char_nonempty_aux_h_exp}
\label{BoolFourier.expectation_char_nonempty_aux_h_exp}
\leanok
\uses{BoolFourier.BoolToPM1, BoolFourier.char_S, BoolFourier.expectation, BoolFourier.expectation_factorizes}
\difficulty{5}
Let $n$ be a natural number and let $S \subseteq [n]$ be a nonempty set of coordinates.
Then the uniform expectation of the Walsh character $\chi_S$ over the Boolean hypercube
factorises as a product over the coordinates in $S$:
\[
  \mathbb{E}[\chi_S]
  \;=\; \prod_{j \in S}
\frac{\mathrm{BoolToPM1}(\mathtt{false}) + \mathrm{BoolToPM1}(\mathtt{true})}{2},
\]
where $\mathrm{BoolToPM1}$ sends $\mathtt{false}$ to $1$ and $\mathtt{true}$ to $-1$.
\end{lemma}

\begin{lemma}[$L^2$ norm squared as a self inner product]
\lean{BoolFourier.parseval_identity_aux_hL2}
\label{BoolFourier.parseval_identity_aux_hL2}
\leanok
\uses{BoolFourier.BoolFun, BoolFourier.L2_norm_sq, BoolFourier.expectation_eq_expect, BooleanAnalysis.innerProduct}
\difficulty{3}
Let $f : \{0,1\}^n \to \bbr$ be a real-valued function on the Boolean hypercube, and
equip such functions with the uniform-measure inner product $\langle g, h \rangle =
2^{-n}\sum_{x \in \{0,1\}^n} g(x)\,h(x)$. Then the squared $L^2$ norm of $f$ equals the
inner product of $f$ with itself:
\[
  \|f\|_2^2 \;=\; \langle f, f \rangle.
\]
\end{lemma}

\begin{lemma}[Interchange of summation in the convolution coefficient]
\lean{BoolFourier.fourier_coeff_convolution_aux_h_fubini}
\label{BoolFourier.fourier_coeff_convolution_aux_h_fubini}
\leanok
\uses{BoolFourier.BoolFun, BoolFourier.char_S, BoolFourier.xor_vec}
\difficulty{2}
Let $n$ be a natural number, let $f, g : \{0,1\}^n \to \bbr$ be Boolean functions, and
let $S \subseteq [n]$ with associated Walsh character $\chi_S(x) = \prod_{i \in
S}(-1)^{x_i}$. Writing $x \oplus y$ for the componentwise XOR of $x, y \in \{0,1\}^n$,
the following identity holds:
\[
\sum_{x \in \{0,1\}^n} \Bigl(\sum_{y \in \{0,1\}^n} f(y)\,g(x \oplus y)\Bigr)\,\chi_S(x)
  \;=\;
  \sum_{y \in \{0,1\}^n} f(y) \sum_{x \in \{0,1\}^n} g(x \oplus y)\,\chi_S(x).
\]
\end{lemma}

\begin{lemma}[Multiplicative splitting of a Walsh character along XOR]
\lean{BoolFourier.fourier_coeff_convolution_aux_h_split}
\label{BoolFourier.fourier_coeff_convolution_aux_h_split}
\leanok
\uses{BoolFourier.BoolFun, BoolFourier.char_S, BoolFourier.hypercube, BoolFourier.xor_vec, BooleanAnalysis.boolToSign, BooleanAnalysis.chiS}
\difficulty{4}
Let $S \subseteq [n]$, let $f, g : \{0,1\}^n \to \bbr$ be Boolean functions, and suppose
the interchange-of-summation identity
\[
  \sum_{x} \Big(\sum_{y} f(y)\, g(x \oplus y)\Big)\,\chi_S(x)
  \;=\; \sum_{y} f(y) \sum_{x} g(x \oplus y)\,\chi_S(x)
\]
holds, where all sums range over $\{0,1\}^n$ and $\oplus$ is componentwise XOR. Then for
all $x, y \in \{0,1\}^n$,
\[
  \chi_S(x) \;=\; \chi_S(x \oplus y)\,\chi_S(y),
\]
where $\chi_S$ is the Walsh--Fourier character of $S$; that is, $\chi_S$ is a character
of the group $(\bbf_2)^n$.
\end{lemma}

\begin{lemma}[Shift rule for character-weighted sums]
\lean{BoolFourier.fourier_coeff_convolution_aux_h_char}
\label{BoolFourier.fourier_coeff_convolution_aux_h_char}
\leanok
\uses{BoolFourier.BoolFun, BoolFourier.char_S, BoolFourier.fourier_coeff_convolution_aux_h_split, BoolFourier.hypercube, BoolFourier.xor_vec}
\difficulty{5}
Let $f, g : \{0,1\}^n \to \bbr$ be Boolean functions, let $S \subseteq [n]$, and write
$\chi_S(x) = \prod_{i \in S}(-1)^{x_i}$ for the associated Walsh character, where
$\oplus$ denotes componentwise XOR. Suppose the order of summation may be interchanged
in the double sum
\[
  \sum_{x}\Big(\sum_{y} f(y)\,g(x \oplus y)\Big)\chi_S(x)
  \;=\;
  \sum_{y} f(y)\sum_{x} g(x \oplus y)\,\chi_S(x).
\]
Then for every $y \in \{0,1\}^n$, translating the argument of $g$ by $y$ multiplies the
character-weighted sum by the factor $\chi_S(y)$:
\[
  \sum_{x} g(x \oplus y)\,\chi_S(x)
  \;=\;
  \chi_S(y)\sum_{x} g(x)\,\chi_S(x).
\]
\end{lemma}
